\chapter{流水线并行 PP}
\label{chap:pipeline-parallelism}

\section{Pipeline Parallelism 的动机}
\label{sec:pipeline-parallelism-motivation}

上一章我们讨论了张量并行。张量并行把一个 Transformer block 内部的 Linear、Attention heads、FFN 中间维度切到多张 GPU 上，使多张 GPU 共同完成一层内部的矩阵计算。但张量并行也有边界：当模型层数很多、总参数量巨大时，即使每一层内部已经被切分，单个 GPU 仍可能无法承担全部层的参数、优化器状态和激活。

这时需要另一类并行方式：\textbf{流水线并行 Pipeline Parallelism, PP}。它的核心思想是：

[
\text{把模型按层切分到不同设备上，让不同 micro-batch 像流水线一样穿过这些 stage。}
]

如果说数据并行是“复制模型、切分数据”，张量并行是“切分层内矩阵”，那么流水线并行就是“切分模型层”。它特别适用于深层 Transformer，因为 Transformer 是由许多相似 block 堆叠而成的结构：

[
\mat{X}_{\ell+1}
================

\operatorname{DecoderBlock}*{\ell}(\mat{X}*{\ell}),
\quad
\ell=0,1,\ldots,L-1.
]

当 $L$ 很大时，可以把这些 block 分给多个 pipeline stages。例如 $L=80$，pipeline size 为 $P=8$，每个 stage 大约负责 $10$ 层。

\subsection{层级切分}
\label{subsec:layer-partitioning}

流水线并行最常见的切分方式是按层切分。设模型有 $L$ 层，pipeline stage 数为 $P$。若均匀切分，每个 stage 拥有：
[
\frac{L}{P}
]
层。第 $p$ 个 stage 负责的层集合可表示为：
[
\mathcal{L}_p
=============

\left{
\ell
\mid
p\frac{L}{P}
\le
\ell
<
(p+1)\frac{L}{P}
\right}.
]

如果 $L$ 不能被 $P$ 整除，则需要不均匀切分。实际系统中，层切分并不总是简单平均，因为不同层的计算量可能不同。例如：

\begin{itemize}
\item embedding 层和 lm head 参数量大，但计算形态不同；
\item 某些模型的前几层或后几层结构特殊；
\item MoE 层与 dense 层成本不同；
\item 使用 recomputation、sequence parallel 或 tensor parallel 后，不同 stage 显存和时间可能不均；
\item 最后一层可能包含 final norm 和 vocab parallel cross entropy。
\end{itemize}

流水线并行要求各 stage 时间尽量均衡。如果某个 stage 比其他 stage 慢，它就会成为整个流水线的瓶颈。设第 $p$ 个 stage 的单个 micro-batch 处理时间为 $t_p$，则流水线稳定阶段的节拍约由最慢 stage 决定：

[
t_{\mathrm{cycle}}
==================

\max_{p} t_p.
]

因此，Pipeline Parallelism 的第一项工程任务就是\textbf{合理切分层，使每个 stage 的计算时间和显存压力尽量接近}。

\subsection{stage}
\label{subsec:pipeline-stage}

Pipeline stage 是流水线并行中的基本计算单元。一个 stage 通常运行在一张 GPU 或一组 GPU 上。如果结合张量并行，一个 pipeline stage 可能不是单张 GPU，而是一个 TP group。例如：

[
\text{PP stage}
===============

\text{一组 TP ranks 共同负责若干层}.
]

在纯流水线并行中，stage 之间传递的是激活张量。前向传播时：

[
\mat{A}_{p+1}
=============

f_p(\mat{A}_p),
]
其中 $f_p$ 表示第 $p$ 个 stage 负责的若干层。第 $p$ 个 stage 计算完成后，将输出激活发送给第 $p+1$ 个 stage。

反向传播时，梯度沿相反方向传播。若第 $p+1$ 个 stage 传回激活梯度：
[
\frac{\partial \mathcal{L}}{\partial \mat{A}*{p+1}},
]
则第 $p$ 个 stage 根据本地保存的前向激活和参数，计算：
[
\frac{\partial \mathcal{L}}{\partial \mat{A}*{p}}
]
并发送给前一个 stage。

因此，流水线并行引入两类点对点通信：

\begin{itemize}
\item forward activation send/recv；
\item backward activation gradient send/recv。
\end{itemize}

这与数据并行和张量并行不同。DDP 主要是 collective all-reduce；TP 常用 all-reduce、all-gather、reduce-scatter；PP 主要是相邻 stage 之间的 point-to-point communication。

\begin{figure}[H]
\centering
\begin{tikzpicture}[
font=\small,
stage/.style={draw, rounded corners=4pt, minimum width=2.1cm, minimum height=1.0cm, align=center, fill=blue!6},
act/.style={draw, rounded corners=2pt, minimum width=1.0cm, minimum height=0.45cm, align=center, fill=orange!14},
arrow/.style={-{Latex[length=2.3mm]}, thick},
back/.style={-{Latex[length=2.3mm]}, thick, red!70}
]

```
\node[font=\bfseries] at (0,3.1) {模型按层切分为多个 Pipeline Stages};

\node[stage] (s0) at (-4.8,0.9) {Stage 0\\Layers 0--9};
\node[stage] (s1) at (-1.6,0.9) {Stage 1\\Layers 10--19};
\node[stage] (s2) at (1.6,0.9) {Stage 2\\Layers 20--29};
\node[stage] (s3) at (4.8,0.9) {Stage 3\\Layers 30--39};

\node[act] (a0) at (-6.4,0.9) {$A_0$};
\node[act] (a1) at (-3.2,0.9) {$A_1$};
\node[act] (a2) at (0.0,0.9) {$A_2$};
\node[act] (a3) at (3.2,0.9) {$A_3$};
\node[act] (a4) at (6.4,0.9) {$A_4$};

\draw[arrow] (a0) -- (s0);
\draw[arrow] (s0) -- (a1);
\draw[arrow] (a1) -- (s1);
\draw[arrow] (s1) -- (a2);
\draw[arrow] (a2) -- (s2);
\draw[arrow] (s2) -- (a3);
\draw[arrow] (a3) -- (s3);
\draw[arrow] (s3) -- (a4);

\draw[back] (a4.south) .. controls (5.6,-0.6) and (4.8,-0.6) .. (s3.south);
\draw[back] (s3.south) .. controls (4.0,-0.95) and (3.2,-0.95) .. (a3.south);
\draw[back] (a3.south) .. controls (2.4,-1.1) and (1.6,-1.1) .. (s2.south);
\draw[back] (s2.south) .. controls (0.8,-1.25) and (0.0,-1.25) .. (a2.south);
\draw[back] (a2.south) .. controls (-0.8,-1.4) and (-1.6,-1.4) .. (s1.south);
\draw[back] (s1.south) .. controls (-2.4,-1.55) and (-3.2,-1.55) .. (a1.south);
\draw[back] (a1.south) .. controls (-4.0,-1.7) and (-4.8,-1.7) .. (s0.south);

\node[align=left, text width=10.7cm] at (0,-2.55) {
  前向传播沿 stage 从左到右传递激活；反向传播沿相反方向传递激活梯度。
  PP 的通信主要发生在相邻 stage 之间，是点对点 send/recv。
};
```

\end{tikzpicture}
\caption{模型按层切分为多个 pipeline stages}
\label{fig:pipeline-stage-layer-partition}
\end{figure}

\subsection{micro-batch}
\label{subsec:micro-batch}

如果只把一个 batch 直接送入流水线，会发生严重空闲。考虑 $P=4$ 个 stages，一个 batch 需要依次经过 Stage 0、Stage 1、Stage 2、Stage 3。若没有切分 micro-batch，则同一时刻只有一个 stage 在工作，其他 stage 都在等待。

为了解决这个问题，流水线并行会把一个 mini-batch 切成多个 micro-batches。设一个训练 step 的 batch 被切成 $M$ 个 micro-batches：

[
\mathcal{B}
===========

{\mathcal{B}_0,\mathcal{B}*1,\ldots,\mathcal{B}*{M-1}}.
]

每个 micro-batch 都会依次经过所有 stages。这样，当 Stage 1 正在处理 micro-batch 0 时，Stage 0 可以处理 micro-batch 1；当 Stage 2 处理 micro-batch 0 时，Stage 1 处理 micro-batch 1，Stage 0 处理 micro-batch 2。多个 micro-batches 填满流水线后，各 stage 可以同时工作。

Global batch size 与 pipeline micro-batch 的关系为：

[
B_{\mathrm{global}}
===================

B_{\mathrm{micro}}
\cdot
M
\cdot
P_{\mathrm{DP}},
]
其中 $B_{\mathrm{micro}}$ 是每个 micro-batch 的 batch size，$M$ 是每个 step 的 micro-batch 数，也常被称为 gradient accumulation steps，$P_{\mathrm{DP}}$ 是数据并行副本数。

需要注意，流水线并行中的 micro-batch 与数据并行中的 local batch、梯度累积密切相关。一个 optimizer step 通常要等待所有 micro-batches 的 forward 和 backward 都完成，然后聚合梯度并更新参数。

\subsection{pipeline bubble}
\label{subsec:pipeline-bubble}

流水线并行的核心问题是 pipeline bubble，即流水线无法完全填满时产生的空闲时间。即使使用 micro-batch，开始阶段需要 warmup，结束阶段需要 drain/cooldown。

以 $P$ 个 stages、$M$ 个 micro-batches 为例。在最简单的流水线中，至少需要 $P-1$ 个时间片填满流水线，也需要 $P-1$ 个时间片排空流水线。因此 bubble 与 $P$ 有关，而可用计算量与 $M$ 有关。粗略地，micro-batch 数越多，bubble 占比越低。

对于简单 GPipe all-forward-all-backward 调度，流水线效率可粗略写作：

[
\eta_{\mathrm{pipe}}
\approx
\frac{M}{M+P-1}.
]

这个公式表达了一个重要直觉：

[
M\gg P
\quad
\Rightarrow
\quad
\eta_{\mathrm{pipe}}\rightarrow 1.
]

也就是说，micro-batch 数越多，流水线越容易被填满，bubble 占比越低。但 $M$ 变大也会带来代价：

\begin{itemize}
\item 每个 optimizer step 的 micro-batch 数增加，step 内调度更复杂；
\item GPipe 中需要保存更多 micro-batch 的激活，显存上升；
\item micro-batch 太小会导致 GEMM shape 变小，计算效率下降；
\item pipeline flush 时间和通信调度变复杂；
\item 与数据并行、张量并行组合时，global batch size 可能过大。
\end{itemize}

\begin{figure}[H]
\centering
\begin{tikzpicture}[
font=\small,
cell/.style={draw, minimum width=0.7cm, minimum height=0.45cm, align=center},
work/.style={fill=blue!14},
bubble/.style={fill=gray!18},
axis/.style={-{Latex[length=2.2mm]}, thick}
]

```
\node[font=\bfseries] at (0,3.1) {Pipeline Bubble：warmup 与 cooldown 中的空闲};

\foreach \t in {0,...,7} {
  \node[font=\scriptsize] at ({-3.5+0.75*\t},2.35) {t\t};
}

\foreach \s/\y in {0/1.7,1/1.15,2/0.6,3/0.05} {
  \node[font=\scriptsize] at (-4.35,\y) {Stage \s};
}

% Stage 0
\foreach \t/\m in {0/0,1/1,2/2,3/3} {
  \node[cell, work] at ({-3.5+0.75*\t},1.7) {$F\m$};
}
\foreach \t in {4,...,7} {
  \node[cell, bubble] at ({-3.5+0.75*\t},1.7) {};
}

% Stage 1
\node[cell, bubble] at (-3.5,1.15) {};
\foreach \t/\m in {1/0,2/1,3/2,4/3} {
  \node[cell, work] at ({-3.5+0.75*\t},1.15) {$F\m$};
}
\foreach \t in {5,...,7} {
  \node[cell, bubble] at ({-3.5+0.75*\t},1.15) {};
}

% Stage 2
\foreach \t in {0,1} {
  \node[cell, bubble] at ({-3.5+0.75*\t},0.6) {};
}
\foreach \t/\m in {2/0,3/1,4/2,5/3} {
  \node[cell, work] at ({-3.5+0.75*\t},0.6) {$F\m$};
}
\foreach \t in {6,7} {
  \node[cell, bubble] at ({-3.5+0.75*\t},0.6) {};
}

% Stage 3
\foreach \t in {0,1,2} {
  \node[cell, bubble] at ({-3.5+0.75*\t},0.05) {};
}
\foreach \t/\m in {3/0,4/1,5/2,6/3} {
  \node[cell, work] at ({-3.5+0.75*\t},0.05) {$F\m$};
}
\node[cell, bubble] at ({-3.5+0.75*7},0.05) {};

\node[draw, rounded corners=2pt, fill=blue!14, minimum width=0.55cm, minimum height=0.35cm] at (3.6,1.3) {};
\node[align=left] at (4.8,1.3) {有效计算};
\node[draw, rounded corners=2pt, fill=gray!18, minimum width=0.55cm, minimum height=0.35cm] at (3.6,0.7) {};
\node[align=left] at (4.8,0.7) {bubble 空闲};

\node[align=left, text width=10.3cm] at (0,-1.45) {
  流水线开始时，后面的 stage 还没有收到输入；结束时，前面的 stage 已经没有新 micro-batch 可处理。
  这些空闲时间就是 pipeline bubble。
};
```

\end{tikzpicture}
\caption{Pipeline warmup/cooldown 造成的 bubble}
\label{fig:pipeline-bubble}
\end{figure}

Pipeline Parallelism 的本质就是在模型切分带来的显存收益与 bubble、通信、调度复杂度之间做权衡。

\section{GPipe}
\label{sec:gpipe}

GPipe 是早期经典的流水线并行调度方式。它的基本策略是：对所有 micro-batches 先执行 forward，等所有 forward 完成后，再按反向顺序执行 backward。这种方式也被称为 all-forward-all-backward。

\subsection{micro-batch 调度}
\label{subsec:gpipe-microbatch-scheduling}

假设有 $P=4$ 个 pipeline stages，$M=4$ 个 micro-batches。GPipe forward 阶段中，micro-batch 0 先进入 Stage 0，下一时间片进入 Stage 1，同时 micro-batch 1 进入 Stage 0。随着时间推进，所有 micro-batches 依次穿过所有 stages。

Forward 完成后，最后一个 stage 开始对最后一个 micro-batch 做 backward，然后梯度逐 stage 向前传播。Backward 阶段与 forward 类似，但方向相反。

GPipe 的优点是概念简单、易于实现。缺点是需要在 backward 前保存所有 micro-batches 的前向激活，因此显存压力较大。

\subsection{all-forward-all-backward}
\label{subsec:all-forward-all-backward}

GPipe 的调度可以写作：

[
F_{0,0},F_{0,1},\ldots,F_{0,P-1},
]
表示 micro-batch 0 从 Stage 0 到 Stage $P-1$ 的 forward；对所有 micro-batches 重复。之后再执行 backward：

[
B_{M-1,P-1},B_{M-1,P-2},\ldots,B_{M-1,0}.
]

更直观地看，它先填充并完成所有 forward，然后再排空 backward。这意味着在 forward 结束时，各 stage 保存了大量 micro-batch 激活，直到对应 backward 使用完毕。

\begin{figure}[H]
\centering
\begin{tikzpicture}[
font=\scriptsize,
cell/.style={draw, minimum width=0.62cm, minimum height=0.42cm, align=center},
fwd/.style={fill=blue!14},
bwd/.style={fill=red!14},
bubble/.style={fill=gray!18}
]

```
\node[font=\bfseries\small] at (0,3.0) {GPipe All-Forward-All-Backward 时间线};

\foreach \t in {0,...,13} {
  \node at ({-4.6+0.65*\t},2.35) {\t};
}

\foreach \s/\y in {0/1.75,1/1.2,2/0.65,3/0.1} {
  \node[font=\scriptsize] at (-5.35,\y) {S\s};
}

% initialize all cells grey
\foreach \s/\y in {0/1.75,1/1.2,2/0.65,3/0.1} {
  \foreach \t in {0,...,13} {
    \node[cell, bubble] at ({-4.6+0.65*\t},\y) {};
  }
}

% Forward schedule M=4, P=4
\foreach \m in {0,...,3} {
  \foreach \s in {0,...,3} {
    \pgfmathtruncatemacro{\t}{\m+\s}
    \pgfmathsetmacro{\y}{1.75-0.55*\s}
    \node[cell, fwd] at ({-4.6+0.65*\t},\y) {$F_{\m}$};
  }
}

% Backward schedule after forward drain, start t=7
\foreach \m/\idx in {3/0,2/1,1/2,0/3} {
  \foreach \s in {3,2,1,0} {
    \pgfmathtruncatemacro{\stageoffset}{3-\s}
    \pgfmathtruncatemacro{\t}{7+\idx+\stageoffset}
    \pgfmathsetmacro{\y}{1.75-0.55*\s}
    \node[cell, bwd] at ({-4.6+0.65*\t},\y) {$B_{\m}$};
  }
}

\node[draw, rounded corners=2pt, fill=blue!14, minimum width=0.5cm, minimum height=0.32cm] at (3.9,1.5) {};
\node[anchor=west] at (4.3,1.5) {Forward};
\node[draw, rounded corners=2pt, fill=red!14, minimum width=0.5cm, minimum height=0.32cm] at (3.9,1.0) {};
\node[anchor=west] at (4.3,1.0) {Backward};
\node[draw, rounded corners=2pt, fill=gray!18, minimum width=0.5cm, minimum height=0.32cm] at (3.9,0.5) {};
\node[anchor=west] at (4.3,0.5) {Bubble};

\node[align=left, text width=10.4cm, font=\small] at (0,-1.35) {
  GPipe 先完成所有 micro-batch 的 forward，再执行所有 backward。
  调度简单，但 forward 激活需要保存到 backward 阶段，显存压力较高。
};
```

\end{tikzpicture}
\caption{GPipe 的 all-forward-all-backward 调度}
\label{fig:gpipe-timeline}
\end{figure}

\subsection{bubble 分析}
\label{subsec:gpipe-bubble-analysis}

对于 $P$ 个 stages、$M$ 个 micro-batches，GPipe forward 阶段需要：
[
M+P-1
]
个时间片完成。Backward 阶段也需要：
[
M+P-1
]
个时间片。因此总时间片约为：
[
T_{\mathrm{GPipe}}
==================

2(M+P-1).
]

真正做有效计算的 forward 和 backward 各有 $MP$ 个 stage-microbatch 任务，总有效任务数为：
[
2MP.
]
如果理想情况下每个时间片每个 stage 都能工作，总容量为：
[
2P(M+P-1).
]
因此效率可近似写为：
[
\eta_{\mathrm{GPipe}}
=====================

# \frac{2MP}{2P(M+P-1)}

\frac{M}{M+P-1}.
]

这就是前面提到的 bubble 公式。

当 $M=1$ 时：
[
\eta_{\mathrm{GPipe}}
=====================

\frac{1}{P}.
]
流水线完全没有被填满，效率很低。

当 $M\gg P$ 时：
[
\eta_{\mathrm{GPipe}}
\approx 1.
]
bubble 占比下降。

但实际中不能无限增大 $M$，因为 micro-batch 数越多，GPipe 需要保存的激活越多。若每个 micro-batch 在每个 stage 保存激活大小为 $A_{\mathrm{stage}}$，则每个 stage 需要保存的激活粗略与 $M$ 成正比：
[
M_{\mathrm{act}}^{\mathrm{GPipe}}
\propto
M A_{\mathrm{stage}}.
]

\subsection{激活保存问题}
\label{subsec:gpipe-activation-saving}

GPipe 的显存问题来自 all-forward-all-backward。所有 micro-batches 的 forward 先完成，backward 尚未开始，因此 stage 必须保存每个 micro-batch 的必要激活。

对于某个 stage，若它处理 $L_p$ 层，每层激活大小约为：
[
B_{\mathrm{micro}}S H s_{\mathrm{dtype}},
]
那么保存 $M$ 个 micro-batches 的激活显存大致与：
[
M L_p B_{\mathrm{micro}} S H s_{\mathrm{dtype}}
]
成正比。实际还包括 attention softmax、MLP 中间激活、dropout mask、norm 中间结果等。

为了缓解这个问题，GPipe 通常与 activation checkpointing 搭配。即每个 stage 内部不保存全部激活，而是在 backward 时重算部分 forward。这样可以显著降低显存，但增加计算量。

下面给出一个简化的 GPipe 调度生成示意代码：

\begin{lstlisting}[language=Python,caption={生成 GPipe all-forward-all-backward 调度的简化示意},label={lst:gpipe-schedule-generator}]
def build_gpipe_schedule(num_stages, num_microbatches):
schedule = []

```
# Forward phase
for t in range(num_microbatches + num_stages - 1):
    tasks = []
    for stage in range(num_stages):
        microbatch = t - stage
        if 0 <= microbatch < num_microbatches:
            tasks.append((stage, "F", microbatch))
    schedule.append(tasks)

# Backward phase
for t in range(num_microbatches + num_stages - 1):
    tasks = []
    for stage in reversed(range(num_stages)):
        # Run backward in reverse stage order.
        microbatch = (num_microbatches - 1) - (t - (num_stages - 1 - stage))
        if 0 <= microbatch < num_microbatches:
            tasks.append((stage, "B", microbatch))
    schedule.append(tasks)

return schedule
```

\end{lstlisting}

这个调度易懂，但不是现代大模型训练中最常见的高效方式。为了降低激活显存并提升 overlap，1F1B 调度更常用。

\section{1F1B Pipeline Schedule}
\label{sec:one-forward-one-backward}

1F1B 是 one-forward-one-backward 的缩写。它的核心思想是：流水线填满后，每个 stage 尽量交替执行一个 forward 和一个 backward，而不是先把所有 forward 做完再做所有 backward。

相比 GPipe，1F1B 的主要优点是减少激活驻留时间。一个 micro-batch 的 forward 完成后，不必等所有 micro-batches 都 forward 完成才 backward。随着 backward 更早开始，部分激活可以更早释放。

1F1B 通常分为三个阶段：

\begin{itemize}
\item \textbf{warmup}：前几个 micro-batches 进入流水线，尚无 backward 可执行；
\item \textbf{steady state}：每个 stage 交替执行 forward 和 backward；
\item \textbf{cooldown}：没有新的 forward，只剩 backward 排空流水线。
\end{itemize}

\subsection{warmup}
\label{subsec:onef1b-warmup}

在 1F1B 中，不同 stage 的 warmup 长度不同。靠前的 stage 需要先产生足够多的 forward activation，让后面的 stage 能够开始 backward。靠后的 stage 较早开始 backward。

对于 $P$ 个 stages，第 $p$ 个 stage 的 warmup micro-batch 数通常与它距离最后 stage 的距离有关。越靠前的 stage，需要更多 warmup forward；越靠后的 stage，warmup 更短。

直观地说，Stage 0 必须先不断发送 micro-batches，直到最后 stage 完成某个 micro-batch 的 forward 并开始 backward，反向梯度再逐 stage 传回来。只有此时 Stage 0 才有 backward 可以执行。

\subsection{steady state}
\label{subsec:onef1b-steady-state}

进入 steady state 后，每个 stage 大致执行：

[
F, B, F, B, F, B,\ldots
]

或者根据实现细节交错执行。这样做有两个收益：

\begin{itemize}
\item backward 更早开始，激活更早释放；
\item stage 的计算更均匀，减少 GPipe 中 forward 与 backward 完全分离带来的激活峰值。
\end{itemize}

1F1B 并不会消除 bubble，但它能显著改善显存。对于大模型训练，这通常比 GPipe 更实用。

\subsection{cooldown}
\label{subsec:onef1b-cooldown}

当所有 micro-batches 都已经完成 forward 后，流水线进入 cooldown。此时不再有新的 forward 输入，剩余任务是把所有 micro-batches 的 backward 执行完。靠后的 stage 较早完成，靠前的 stage 最后完成 backward。

Cooldown 与 warmup 一样会产生 bubble，但当 micro-batch 数足够大时，bubble 占比可以被摊薄。

\begin{figure}[H]
\centering
\begin{tikzpicture}[
font=\scriptsize,
cell/.style={draw, minimum width=0.62cm, minimum height=0.42cm, align=center},
fwd/.style={fill=blue!14},
bwd/.style={fill=red!14},
bubble/.style={fill=gray!18}
]

```
\node[font=\bfseries\small] at (0,3.0) {1F1B Pipeline 时间线示意};

\foreach \t in {0,...,13} {
  \node at ({-4.6+0.65*\t},2.35) {\t};
}

\foreach \s/\y in {0/1.75,1/1.2,2/0.65,3/0.1} {
  \node[font=\scriptsize] at (-5.35,\y) {S\s};
}

\foreach \s/\y in {0/1.75,1/1.2,2/0.65,3/0.1} {
  \foreach \t in {0,...,13} {
    \node[cell, bubble] at ({-4.6+0.65*\t},\y) {};
  }
}

% A hand-crafted 1F1B-like timeline for P=4, M=6
% Stage 0
\foreach \t/\op/\m/\style in {
  0/F/0/fwd,1/F/1/fwd,2/F/2/fwd,3/F/3/fwd,4/B/0/bwd,5/F/4/fwd,6/B/1/bwd,7/F/5/fwd,8/B/2/bwd,9/B/3/bwd,10/B/4/bwd,11/B/5/bwd
} {
  \node[cell, \style] at ({-4.6+0.65*\t},1.75) {$\op_{\m}$};
}

% Stage 1
\foreach \t/\op/\m/\style in {
  1/F/0/fwd,2/F/1/fwd,3/F/2/fwd,4/B/0/bwd,5/F/3/fwd,6/B/1/bwd,7/F/4/fwd,8/B/2/bwd,9/F/5/fwd,10/B/3/bwd,11/B/4/bwd,12/B/5/bwd
} {
  \node[cell, \style] at ({-4.6+0.65*\t},1.2) {$\op_{\m}$};
}

% Stage 2
\foreach \t/\op/\m/\style in {
  2/F/0/fwd,3/F/1/fwd,4/B/0/bwd,5/F/2/fwd,6/B/1/bwd,7/F/3/fwd,8/B/2/bwd,9/F/4/fwd,10/B/3/bwd,11/F/5/fwd,12/B/4/bwd,13/B/5/bwd
} {
  \node[cell, \style] at ({-4.6+0.65*\t},0.65) {$\op_{\m}$};
}

% Stage 3
\foreach \t/\op/\m/\style in {
  3/F/0/fwd,4/B/0/bwd,5/F/1/fwd,6/B/1/bwd,7/F/2/fwd,8/B/2/bwd,9/F/3/fwd,10/B/3/bwd,11/F/4/fwd,12/B/4/bwd,13/F/5/fwd
} {
  \node[cell, \style] at ({-4.6+0.65*\t},0.1) {$\op_{\m}$};
}

\node[draw, rounded corners=2pt, fill=blue!14, minimum width=0.5cm, minimum height=0.32cm] at (3.9,1.5) {};
\node[anchor=west] at (4.3,1.5) {Forward};
\node[draw, rounded corners=2pt, fill=red!14, minimum width=0.5cm, minimum height=0.32cm] at (3.9,1.0) {};
\node[anchor=west] at (4.3,1.0) {Backward};
\node[draw, rounded corners=2pt, fill=gray!18, minimum width=0.5cm, minimum height=0.32cm] at (3.9,0.5) {};
\node[anchor=west] at (4.3,0.5) {Bubble};

\node[align=left, text width=10.4cm, font=\small] at (0,-1.4) {
  1F1B 在 warmup 后交替执行 forward 与 backward。
  相比 GPipe，它能让激活更早释放，因此更适合显存紧张的大模型训练。
};
```

\end{tikzpicture}
\caption{1F1B pipeline schedule 的 warmup、steady state 与 cooldown}
\label{fig:onef1b-timeline}
\end{figure}

\subsection{与 GPipe 的对比}
\label{subsec:gpipe-vs-onef1b}

GPipe 和 1F1B 都使用 micro-batch 填充流水线，但它们的显存和调度特征不同。

\begin{table}[H]
\centering
\small
\caption{GPipe 与 1F1B Pipeline Schedule 对比}
\label{tab:gpipe-vs-1f1b}
\begin{tabular}{p{0.2\textwidth}p{0.34\textwidth}p{0.34\textwidth}}
\toprule
\textbf{维度} & \textbf{GPipe} & \textbf{1F1B} \
\midrule
调度方式
& 所有 forward 完成后再执行所有 backward
& warmup 后 forward/backward 交替执行 \
\midrule
实现复杂度
& 较简单
& 较复杂，需要更精细的 send/recv 调度 \
\midrule
激活显存
& 较高，需要保存所有 micro-batch 激活
& 较低，部分 micro-batch 可更早 backward 并释放激活 \
\midrule
pipeline bubble
& 存在，micro-batch 数越多占比越低
& 也存在，但 steady state 更适合长时间运行 \
\midrule
大模型适用性
& 概念清晰，适合教学和部分场景
& 更常用于大规模模型训练 \
\bottomrule
\end{tabular}
\end{table}

1F1B 的一个关键点是梯度累积。由于一个 optimizer step 包含多个 micro-batches，每个 stage 在完成所有 micro-batches 的 backward 之前不能更新参数。否则不同 micro-batch 会看到不同版本的参数，破坏同步语义。

因此，1F1B 中通常存在一个参数版本一致性要求：

[
\theta_{\mathrm{used\ by\ all\ microbatches}}
=============================================

\theta^{(t)}.
]
直到本 step 所有 micro-batches backward 完成后，才执行：

[
\theta^{(t+1)}
==============

\operatorname{OptimizerStep}
\left(
\theta^{(t)},\nabla\theta
\right).
]

\section{Interleaved Pipeline}
\label{sec:interleaved-pipeline}

随着模型越来越大，单纯把连续层分配给每个 physical stage 可能仍然存在 bubble 和负载不均问题。Interleaved Pipeline 使用 virtual pipeline stage，把每个物理设备上的层再切成多个虚拟段，交错执行，以减少 bubble 并改善负载均衡。

\subsection{virtual pipeline stage}
\label{subsec:virtual-pipeline-stage}

假设有 $P$ 个 physical pipeline stages，每个 stage 原本负责连续的一段层。如果引入 virtual pipeline size $V$，则每个 physical stage 被拆成 $V$ 个 virtual stages。总 virtual stages 数为：
[
P_{\mathrm{virtual}}=P\cdot V.
]

例如模型有 $32$ 层，physical PP size 为 $4$，每个 stage 原本负责 $8$ 层：

[
\text{Stage 0}: L0-L7,
\quad
\text{Stage 1}: L8-L15,
]
[
\text{Stage 2}: L16-L23,
\quad
\text{Stage 3}: L24-L31.
]

若 virtual pipeline size 为 $2$，则每个 physical stage 可以拥有两段非连续或交错的层：

[
\text{GPU0}: L0-L3,\ L16-L19,
]
[
\text{GPU1}: L4-L7,\ L20-L23,
]
[
\text{GPU2}: L8-L11,\ L24-L27,
]
[
\text{GPU3}: L12-L15,\ L28-L31.
]

这种布局使同一物理 GPU 在流水线中扮演多个 virtual stage，从而缩短每个 virtual stage 的计算时间，使流水线粒度更细。

\begin{figure}[H]
\centering
\begin{tikzpicture}[
font=\small,
gpu/.style={draw, rounded corners=4pt, minimum width=2.1cm, minimum height=1.35cm, align=center, fill=blue!6},
vstage/.style={draw, rounded corners=2pt, minimum width=1.55cm, minimum height=0.42cm, align=center, fill=orange!14},
arrow/.style={-{Latex[length=2.2mm]}, thick}
]

```
\node[font=\bfseries] at (0,3.2) {Interleaved Pipeline：一个物理 GPU 包含多个 Virtual Stages};

\foreach \i/\x in {0/-4.8,1/-1.6,2/1.6,3/4.8} {
  \node[gpu] (g\i) at (\x,0.9) {GPU \i};
  \node[vstage] (a\i) at (\x,1.25) {Virtual Stage A};
  \node[vstage] (b\i) at (\x,0.55) {Virtual Stage B};
}

\draw[arrow] (a0) -- (a1);
\draw[arrow] (a1) -- (a2);
\draw[arrow] (a2) -- (a3);
\draw[arrow] (a3.south) .. controls (3.6,-0.55) and (-3.6,-0.55) .. (b0.south);
\draw[arrow] (b0) -- (b1);
\draw[arrow] (b1) -- (b2);
\draw[arrow] (b2) -- (b3);

\node[align=left, text width=10.5cm] at (0,-1.65) {
  每个物理 GPU 负责多个 virtual stages。模型层在 virtual pipeline 中更细粒度流动，
  可以降低 bubble，但调度、通信和激活管理更复杂。
};
```

\end{tikzpicture}
\caption{Interleaved Pipeline 中的 virtual pipeline stage}
\label{fig:interleaved-virtual-pipeline-stage}
\end{figure}

\subsection{减少 bubble}
\label{subsec:interleaved-reduce-bubble}

Interleaving 的核心收益是减少 pipeline bubble。直觉上，virtual stage 数增多后，流水线更细，单个 stage 的处理时间更短，填充和排空的粒度更细。对于相同 physical PP size，interleaved schedule 可以提高 pipeline 利用率。

但收益不是免费的。Interleaving 会增加：

\begin{itemize}
\item stage 间激活传递次数；
\item send/recv 调度复杂度；
\item 每个 GPU 上多个 virtual stage 的激活管理；
\item layer 到 device 的映射复杂度；
\item checkpoint 保存和恢复的映射复杂度。
\end{itemize}

如果 virtual stage 过多，单个 virtual stage 的计算太小，通信和调度开销可能抵消 bubble 收益。因此 virtual pipeline size 也需要调参。

\subsection{调度复杂度}
\label{subsec:interleaved-scheduling-complexity}

Interleaved schedule 需要同时考虑：

\begin{itemize}
\item physical stage；
\item virtual stage；
\item micro-batch id；
\item forward/backward 方向；
\item 激活 send/recv；
\item 梯度 send/recv；
\item 与 tensor parallel communication 的重叠；
\item 与 data parallel gradient reduction 的关系。
\end{itemize}

一个训练 step 中，某个 GPU 可能在不同时间片执行不同 virtual stage 的 forward 或 backward。它还要保存每个 micro-batch 对应 virtual stage 的激活。若实现不当，很容易出现死锁、显存泄漏、激活错配或梯度错配。

下面给出一个抽象的 interleaved task 表示：

\begin{lstlisting}[language=Python,caption={Interleaved Pipeline 中任务描述的抽象表示},label={lst:interleaved-task-abstraction}]
from dataclasses import dataclass

@dataclass
class PipelineTask:
physical_stage: int
virtual_stage: int
microbatch: int
direction: str  # "forward" or "backward"

def run_task(task: PipelineTask):
if task.direction == "forward":
recv_activation_if_needed(task)
output = run_forward_layers(task)
send_activation_if_needed(task, output)
save_activation_for_backward(task, output)
else:
recv_grad_if_needed(task)
grad_input = run_backward_layers(task)
send_grad_if_needed(task, grad_input)
release_saved_activation(task)
\end{lstlisting}

真实系统会比这复杂得多。尤其是当 pipeline parallel 与 tensor parallel、data parallel、sequence parallel、activation checkpointing 同时存在时，每个任务可能同时涉及 point-to-point communication 和 collective communication。

\subsection{Megatron PP 实践}
\label{subsec:megatron-pipeline-practice}

Megatron-LM 的 3D 并行通常组合 TP、PP 和 DP。对于一个 rank，它同时属于：

[
\text{tensor model parallel group},
]
[
\text{pipeline model parallel group},
]
[
\text{data parallel group}.
]

Pipeline rank 决定它负责哪些层；Tensor parallel rank 决定它负责这些层中哪些 tensor shards；Data parallel rank 决定它属于哪个 replica。

如果 world size 为：
[
W=P_{\mathrm{TP}}\cdot P_{\mathrm{PP}}\cdot P_{\mathrm{DP}},
]
则总 rank 会被组织成三维网格。后续 3D 并行章节会详细讨论 rank mapping。这里先建立直觉：Pipeline Parallelism 不会替代 TP 或 DP，而是与它们组合使用。

在 Megatron 风格训练中，流水线并行实践需要考虑：

\begin{itemize}
\item 每个 pipeline stage 的层数是否均衡；
\item embedding 和 lm head 放在哪些 stage；
\item 是否使用 interleaved schedule；
\item micro-batch 数是否足够填满流水线；
\item activation checkpointing 的粒度；
\item stage 间通信是否跨节点；
\item 与 TP collective 是否发生资源竞争；
\item 与 DP gradient synchronization 的重叠方式。
\end{itemize}

\section{流水线并行的工程权衡}
\label{sec:pipeline-parallel-engineering-tradeoffs}

\subsection{显存收益}
\label{subsec:pp-memory-benefit}

流水线并行按层切分模型，因此每个 stage 只保存部分层的参数、梯度和优化器状态。若模型层均匀切分到 $P_{\mathrm{PP}}$ 个 stages，理想情况下每个 stage 的层相关参数显存约下降为：

[
\frac{1}{P_{\mathrm{PP}}}.
]

但并非所有显存都严格按 PP size 缩小。例如：

\begin{itemize}
\item embedding 和 lm head 可能集中在首尾 stage；
\item pipeline boundary 需要保存发送和接收 buffer；
\item micro-batch 数增加会影响激活数量；
\item GPipe 会保存大量 forward 激活；
\item 1F1B 降低激活驻留，但仍需要保存未 backward 的 micro-batch；
\item checkpointing 会用额外计算换显存。
\end{itemize}

可以粗略写作：

[
M_{\mathrm{per\ stage}}
\approx
\frac{M_{\mathrm{layer\ states}}}{P_{\mathrm{PP}}}
+
M_{\mathrm{stage\ activations}}
+
M_{\mathrm{pipeline\ buffers}}
+
M_{\mathrm{imbalanced}}.
]

其中 $M_{\mathrm{imbalanced}}$ 表示由不均匀层切分、embedding/lm head 或特殊层造成的额外显存不平衡。

\subsection{通信成本}
\label{subsec:pp-communication-cost}

PP 的通信主要是 stage boundary 的 activation 和 activation gradient。设 stage 间传递的激活 shape 为：
[
\mathbb{R}^{B_{\mathrm{micro}}\times S\times H}.
]
若使用 BF16，每个元素 2 bytes，则一次 forward 边界通信大小约为：
[
B_{\mathrm{micro}} S H \cdot 2\ \mathrm{bytes}.
]
Backward 还需要传回同样大小的梯度。因此每个 micro-batch 每条边界的双向通信约为：
[
2B_{\mathrm{micro}} S H s_{\mathrm{dtype}}.
]

对于 $P_{\mathrm{PP}}$ 个 stages，有 $P_{\mathrm{PP}}-1$ 条边界。每个 step 有 $M$ 个 micro-batches。因此 stage boundary 通信总量与：
[
M(P_{\mathrm{PP}}-1)B_{\mathrm{micro}}S Hs_{\mathrm{dtype}}
]
成正比。

与 TP 不同，PP 通信通常是相邻 stages 的 point-to-point，通信量与 hidden activation 大小有关，而不是与参数量直接相关。长序列和大 hidden 会增加 PP 通信压力。

\subsection{负载均衡}
\label{subsec:pp-load-balance}

Pipeline 的吞吐由最慢 stage 决定。若某个 stage 时间为 $t_{\max}$，其他 stages 更快，则更快的 stages 会等待它。负载不均来自：

\begin{itemize}
\item 层数切分不均；
\item embedding/lm head 集中在首尾 stage；
\item MoE 层造成某些 stage 更慢；
\item 激活 checkpointing 只在部分 stage 开启；
\item 硬件性能差异或节点温度降频；
\item stage 间通信跨越不同带宽链路。
\end{itemize}

负载均衡不仅要看理论 FLOPs，还要看 profiler 中的实际 stage time。某个 stage 即使 FLOPs 不高，也可能因为通信、内存带宽或小 kernel 过多而变慢。

\subsection{调试难度}
\label{subsec:pp-debugging-difficulty}

Pipeline Parallelism 的调试比 DDP 和单纯 TP 更复杂。常见问题包括：

\begin{itemize}
\item send/recv 不匹配导致死锁；
\item 某个 rank 提前退出导致其他 rank hang；
\item micro-batch id 与 activation buffer 错配；
\item backward gradient 发送给错误 stage；
\item 不同 stage 参数版本不一致；
\item pipeline flush 时遗漏某些 backward；
\item checkpoint 保存的 layer 到 stage 映射错误；
\item stage 负载不均导致吞吐低；
\item 日志只在 rank 0 打印，难以看到中间 stage 错误。
\end{itemize}

因此，工程实践中通常需要：

\begin{itemize}
\item 为每个 rank/stage 打印可控调试日志；
\item 在小模型、小 micro-batch 上验证 schedule；
\item 使用张量 shape assertion；
\item 在通信前后记录 micro-batch id；
\item 使用 profiler 查看每个 stage timeline；
\item 将调度逻辑与模型层逻辑解耦；
\item 提供 deterministic 的 pipeline schedule 测试。
\end{itemize}

\section{一个极简流水线模拟器}
\label{sec:minimal-pipeline-simulator}

为了帮助理解 pipeline bubble，可以写一个极简模拟器。它不执行真实模型，只生成每个时间片每个 stage 的任务。

\begin{lstlisting}[language=Python,caption={极简 GPipe 时间线模拟器},label={lst:minimal-gpipe-simulator}]
def gpipe_timeline(num_stages, num_microbatches):
total_forward_steps = num_microbatches + num_stages - 1
total_backward_steps = num_microbatches + num_stages - 1

```
timeline = []

# Forward
for t in range(total_forward_steps):
    row = []
    for s in range(num_stages):
        m = t - s
        if 0 <= m < num_microbatches:
            row.append(f"F{m}")
        else:
            row.append("--")
    timeline.append(row)

# Backward
for t in range(total_backward_steps):
    row = []
    for s in range(num_stages):
        # backward flows from last stage to first
        m = num_microbatches - 1 - (t - (num_stages - 1 - s))
        if 0 <= m < num_microbatches:
            row.append(f"B{m}")
        else:
            row.append("--")
    timeline.append(row)

return timeline
```

def print_timeline(timeline):
num_stages = len(timeline[0])
print("time | " + " ".join([f"S{s:>3}" for s in range(num_stages)]))
print("-" * (7 + 5 * num_stages))
for t, row in enumerate(timeline):
print(f"{t:>4} | " + " ".join([f"{x:>4}" for x in row]))

timeline = gpipe_timeline(num_stages=4, num_microbatches=6)
print_timeline(timeline)
\end{lstlisting}

对于 1F1B，调度生成会更复杂，需要根据 stage、micro-batch、warmup 数、剩余 backward 数决定下一步执行 forward 还是 backward。真实框架通常不会用一个简单循环描述完整调度，而是使用状态机或预生成 schedule table。

下面给出一个更抽象的 1F1B 状态机伪代码：

\begin{lstlisting}[language=Python,caption={1F1B 调度的状态机式伪代码},label={lst:onef1b-state-machine}]
class PipelineStageState:
def **init**(self, stage_id, num_stages, num_microbatches):
self.stage_id = stage_id
self.num_stages = num_stages
self.num_microbatches = num_microbatches

```
    self.next_forward = 0
    self.next_backward = 0
    self.saved_activations = {}

def has_forward_to_run(self):
    return self.next_forward < self.num_microbatches

def has_backward_to_run(self):
    # In a real implementation, this depends on receiving grad from next stage.
    return self.next_backward < self.next_forward - (self.num_stages - 1 - self.stage_id)

def step(self):
    if self.in_warmup_phase():
        return self.run_forward()
    elif self.has_backward_to_run():
        # In steady state, prefer backward to release activation.
        return self.run_backward()
    elif self.has_forward_to_run():
        return self.run_forward()
    else:
        return self.idle_or_cooldown()

def run_forward(self):
    mb = self.next_forward
    self.next_forward += 1
    return ("F", mb)

def run_backward(self):
    mb = self.next_backward
    self.next_backward += 1
    return ("B", mb)
```

\end{lstlisting}

这段伪代码不是可直接用于训练的实现，但它说明了 1F1B 的核心：调度器要知道每个 stage 当前能否 forward、能否 backward、是否已经收到必要输入或梯度，并尽可能优先执行能释放激活的 backward。

\section{本章小结}
\label{sec:chapter09-summary}

本章系统介绍了流水线并行 Pipeline Parallelism 的动机、基本概念、GPipe、1F1B 和 Interleaved Pipeline。

\begin{itemize}
\item \textbf{流水线并行的核心思想} 是把模型按层切成多个 pipeline stages，让不同 micro-batch 像流水线一样穿过这些 stages。
\item \textbf{Stage} 是流水线中的基本计算单元，负责一段连续或虚拟切分的模型层；stage 间传递 forward activation 和 backward activation gradient。
\item \textbf{Micro-batch} 用于填满流水线。Global batch size 满足 $B_{\mathrm{global}}=B_{\mathrm{micro}}\cdot M\cdot P_{\mathrm{DP}}$。
\item \textbf{Pipeline bubble} 是 warmup 和 cooldown 中的空闲时间。简单情况下，流水线效率近似为 $\eta_{\mathrm{pipe}}\approx \frac{M}{M+P-1}$。
\item \textbf{GPipe} 使用 all-forward-all-backward 调度，概念简单，但需要保存所有 micro-batch 的 forward 激活，显存压力较大。
\item \textbf{1F1B} 在 warmup 后交替执行 forward 和 backward，能让激活更早释放，更适合大模型训练。
\item \textbf{Interleaved Pipeline} 引入 virtual pipeline stage，让每个物理 GPU 承担多个虚拟段，以减少 bubble，但调度和通信更复杂。
\item \textbf{PP 的显存收益} 来自按层切分模型状态，但激活、pipeline buffer、embedding/lm head 和负载不均会影响真实显存。
\item \textbf{PP 的通信成本} 主要来自相邻 stage 之间传递激活和激活梯度，通信量与 micro-batch、sequence length 和 hidden size 相关。
\item \textbf{PP 的工程难点} 包括 stage 负载均衡、send/recv 死锁、micro-batch 调度、参数版本一致性、checkpoint 映射和 profiler 分析。
\end{itemize}

下一章将进入 3D 并行策略。我们会把数据并行、张量并行和流水线并行组合起来，分析 DP、TP、PP 三者如何正交组织 rank group，如何根据 GPU 数量反推并行配置，以及如何在计算、显存和通信之间做系统级权衡。
